\section{Introduction}\label{introduction}

The persistent lack of synchronization between demand and capacity
continues to present significant challenges, even for modern
manufacturing systems (Guo \emph{et al.}, 2024). This mismatch often
arises from the inherent variability within these systems, driven by
factors such as market fluctuations, resource availability, and supply
chain disruptions (Hopp and Spearman, 2021). Effective production
scheduling plays a pivotal role in synchronizing demand with capacity by
defining what to produce, how much to produce, where to produce, and in
what sequence (\colorbox{yellow!40}{cf., Reference ≥ 2021}). Traditionally, production
scheduling has relied on various approaches, ranging from simple
dispatching rules (\colorbox{yellow!40}{e.g., Reference}) to more sophisticated
algorithms (\colorbox{yellow!40}{e.g., Reference ≥ 2021}). Recent advances in the field
have increasingly incorporated meta-heuristics to achieve optimal or
near-optimal solutions (\colorbox{yellow!40}{Reference ≥ 2021}). However, using
meta-heuristics in production scheduling presents challenges,
particularly in dealing with the uncertainty and dynamic conditions
typical of manufacturing environments (\colorbox{yellow!40}{Reference}). Meta-heuristics
such as genetic algorithms, simulated annealing, or tabu search are
often designed for deterministic problems, where environmental
conditions and decision outcomes are fixed (\colorbox{yellow!40}{cf. Reference}). In
contrast, stochastic optimization involves uncertainty and probabilistic
outcomes, which meta-heuristics may struggle to accommodate, often
leading to suboptimal solutions or long computation times
(\colorbox{yellow!40}{Reference}). These algorithms typically rely on heuristics that are
not well-suited to random or fluctuating environments, especially when
applied to large-scale or complex scheduling problems (Márquez and
Ribeiro, 2022). Additionally, meta-heuristics may fail to fully
capitalize on the stochastic nature of the problem, lacking the ability
to learn from past experiences or adjust strategies based on observed
randomness (\colorbox{yellow!40}{Reference}). As a result, this can lead to inefficient
exploration of the solution space and a higher likelihood of converging
to a local optimum rather than a global one (\colorbox{yellow!40}{Reference}).
Furthermore, meta-heuristic algorithms are typically tailored to
specific problems, making them less versatile and prone to poor
generalization (\colorbox{yellow!40}{Reference}). To overcome these challenges,
reinforcement learning (RL) has emerged as a promising alternative for
production scheduling (\colorbox{yellow!40}{Reference ≥ 2021}). RL is a branch of machine
learning where an agent interacts with its environment to maximize
long-term rewards (Sutton and Barto, 2018). Its capacity to make
sequential decisions and adapt to stochastic environments makes RL
especially suitable for production scheduling applications (\colorbox{yellow!40}{cf., Li
et al., 2022; del Real Torres et al., 2022; Wang et al., 2022}).

RL has proven effective in various manufacturing contexts. For instance,
Tang and Salonitis (\colorbox{yellow!40}{2021}) leveraged deep RL with discrete-event
simulation to optimize scheduling in reconfigurable systems, surpassing
conventional dispatching rules. Chen et al. (\colorbox{yellow!40}{2022}) addressed the
flexible job shop scheduling problem with resource preemption using a
deep RL-based algorithm, outperforming traditional meta-heuristics. Liu
et al. (\colorbox{yellow!40}{2022}) applied a Deep Q-Network algorithm for real-time
scheduling in dynamic flexible job shop manufacturing. Gerpott et al.
(\colorbox{yellow!40}{2022}) used the Actor-Critic algorithm to minimize tardiness and
makespan in a two-stage hybrid flow shop. Beyond standalone
applications, RL has been integrated with some operations management
(OM) theories. For example, Silva and Azevedo (\colorbox{yellow!40}{2019}) applied RL to
optimize work-in-process (WIP) caps in CONWIP and push systems,
achieving a substantial reduction in WIP without sacrificing production
volume. Xanthopoulos et al. (\colorbox{yellow!40}{2019}) demonstrated RL's superiority in
deriving optimal control policies for CONWIP-type systems. Paraschos et
al. (\colorbox{yellow!40}{2023}) used RL to optimize a three-stage production line,
promoting Lean practices. Lahrichi et al. (\colorbox{yellow!40}{2023}) used SARSA
algorithm to dynamically parameterize demand-driven material
requirements planning, reducing inventory and preventing stockouts.

One OM theory that has yet to be explored with RL in production
scheduling is the Theory of Constraints (TOC). Like RL, which addresses
decision-making in stochastic environments, TOC also does so, but with a
strict focus on the constraint that limits the system's overall
performance, usually protecting it through the use of buffers (Goldratt
and Cox, 1984a). In production scheduling, TOC introduces two prominent
techniques (Goldratt, 2006a; Schragenheim and Ronen, 1990): (i)
drum-buffer-rope (DBR); and (ii) buffer management (BM). When applied
together, these techniques, along with their associated performance
measures --- throughput, inventory, and operating expenses --- enable
the effective management of highly interdependent systems
(\colorbox{yellow!40}{Reference}). They achieve this by protecting capacity-constrained
resources against shortages or blockages through inventory buffers while
synchronizing inflows to the rhythm of the constraint (\colorbox{yellow!40}{Reference}).
TOC-based solutions to scheduling problems have demonstrated their
effectiveness in various contexts. For instance, they have improved
efficiency in production lines within aerospace engineering-to-order
facilities (Telles \emph{et al.}, 2022a) and increased productivity in
small and medium-sized manufacturing enterprises (Buestán Benavides and
Van Landeghem, 2015). Furthermore, these solutions have been
successfully applied to automotive parts manufacturing (Golmohammadi,
2015), LCD panel production (Wu \emph{et al.}, 2010), and engine
cylinder restoration (Ma and Chen, 2009).

Despite the promising synergy between RL and TOC in production
scheduling, their integration remains largely unexplored. This gap
prompts the following research questions: (i) how can RL and TOC be
combined to address production scheduling in stochastic manufacturing
environments? (ii) what practical and theoretical outcomes can result
from this combination? To address these questions, we employed
computational simulation and conducted two experiments within the same
stochastic manufacturing environment. In the first experiment,
production scheduling was performed using a TOC-based solution (i.e.,
DBR and BM). In the second, scheduling was carried out combining RL with
TOC. From a pragmatic perspective, a summative evaluation comparing
these experiments demonstrated that a scheduling solution combining RL
and TOC outperformed the TOC-based solution in terms of both production
volume and inventory levels. Theoretically, a formative evaluation
revealed that this advantage stems from the RL agent's ability to
monitor and adapt to system-wide WIP dynamics, rather than focusing
solely on the constraint. We embedded this adaptive behavior into the
traditional TOC scheduling solution, enhancing its DBR and BM
components, which we believe is the main contribution of this study.

The remainder of this paper is structured as follows. Section 2 provides
the theoretical background on both RL and the TOC in the context of
production scheduling. Section 3 outlines the research design, while
Section 4 details the development of simulation models. Section 5
describes the setup of the simulation experiments, with the
corresponding summative and formative evaluations presented in Section
6. Finally, Section 7 discusses the results, and Section 8
contextualizes the research contributions and limitations, highlighting
opportunities for future studies.
